\scanpage{27}
\mspage{12}

L'extension $\unsure{\Sigma|E}=\Sigma_{0,3}$ de $\Gamma=\Gamma_{03}=\underbrace{\unsure{\mathfrak S|\mathcal S}_{[0,1,\infty]}}_{\text{ou }\mathfrak S_3}\times\mathbb Z/2\mathbb Z$ \inserted{engendré par $\tau:z\mapsto\overline z$} par $\pi=\pi_{0,3}=\{l_0,l_1,l_\infty\mid l_\infty l_1l_0=1\}=\pi_1(\underbrace{\mathbb P^1_{\unsure{\mathbb C|C}}\setminus\{0,1,\infty\}}_{U=U_{0,3}},-\bar j\}$.

Les classes de conjugaison de scindages partiels maximaux \deleted{au-dessus de} correspondent aux points fixes de $\dot\rho,\dot\sigma_0,\dot\sigma_1,\dot\sigma_\infty$ opérant sur $X_{0,3}$, et \inserted{(éventuellt)} aux comp. connexes des ens. d'invariants de $\deleted{\unsure{\rho}}\tau,\dot\sigma_0\tau,\dot\sigma_1\tau,\dot\sigma_\infty\tau$. Les pts fixes de $\rho$ sont $P$ et $\overline P$ dont le stabilisateur est d'ordre 6, \inserted{savoir $\varphi(\mathfrak S_3)$} formé de $1,\dot\rho,\dot\rho^2,\dot\sigma_1\tau,\dot\sigma_2\tau,\dot\sigma_3\tau$.

On a étudié le relèvement correspondant \inserted{à $P$} de $\varphi(\mathfrak S_3)$, on trouve des relèvements ($\pi$-conjugués \inserted{entre eux}) relatifs à $\overline P$ s'obtiennent en \deleted{\unsure{conj}} prenant un élément de $\Sigma$ \deleted{\unsure{q}} qui \deleted{\unsure{soit|sont}} au dessus d'un élément de $\Gamma$ envoyant $P$ en $\overline P$ (savoir $\dot\sigma_0,\dot\sigma_1,\dot\sigma_\infty,\tau,\tau\rho,\tau\rho^2$) et en conjuguant par cet élément

on peut donc prendre $\deleted{\unsure{\sigma_0,\tau}}$ un des éléments $\sigma_0,\sigma_1,\sigma_\infty$ (ou $\varepsilon_0,\varepsilon_1,\varepsilon_\infty$) ou $\tau_0,\tau_1,\tau_\infty$ ou \deleted{\unsure{les}} les produits des précédents par \deleted{\unsure{un}} $\rho$ ou $\rho^2$\ldots.

\deleted{Les} Le \inserted{seul} pt\deleted{s} fixe\deleted{s} de $\sigma_0$ est 2, dont le stabilisateur est le groupe $\simeq\mathbb Z/2\times\mathbb Z/2$ engendré par $\sigma_0$ et $\tau$. On a un relèvt de ce sous-groupe
% archive sheet number: 26
